\chapter{计算材料科学概述}

%二十世纪七十年代以来，随着计算机科学和密度泛函理论\textrm{(Density Functional Theory， DFT)}\cite{PRB136-864_1964,PRA140-1133_1965,Parr-Yang,CR91-651_1991}、分子动力学\textrm{(Molecular Dynamics， MD)}的发展和融合，计算材料科学\textrm{(computational materials science)}与物理、化学、工程力学以及应用数学等诸多基础和应用学科交叉日益紧密，逐渐成为一门新兴的独立学科。\cite{Nat-Mat3-429_2004,App-CataA5-254_2003,JACS125-4306_2003,JCC5-472_2003,Mater_Sci-Tech16-1_2005,Nature392-694_1998}材料计算模拟研究的不断深入，也推动了材料模拟核心计算软件日益成熟。核心计算软件解决的是物质运动基本方程求解和基本物性计算的问题，除此之外，完整的材料模拟过程还包括“计算前处理”的问题建模、计算参数选择到“计算后处理”的数据分析，几乎每一部分都有大量具体的工作有赖人工，频繁的人机交互严重影响了模拟计算流程的顺畅与效率。进入二十一世纪，大规模高性能计算的兴起，逐渐变革了材料计算的研发模式。通过计算流程设计，衔接不同尺度的计算过程，降低人机交互频次，优化模拟流程中的任务调度，实现材料计算流程的自动化与跨尺度物性模拟，形成完善的材料数据库，已经成为材料计算软件发展的新趋势。


计算材料科学是计算机科学与凝聚态物理、理论与计算化学、材料科学及应用数学等诸多基础学科交叉融合中形成的一门新兴的前沿学科。它试图基于基础理论构建的物理模型，通过数值模拟与计算，揭示材料的微观构成、物理性质与宏观使用性能之间的复杂的数据关系%，被视为继理论研究和实验科学之后的“第三种科学方法”
。计算材料科学的核心理念是通过对材料结构-性能之间关系的探索，实现材料制备与加工过程的定量化控制，从而缩短材料研制周期、控制并降低制造成本，提升材料的服役性能，进而为按需设计新材料开辟新的技术路径。%美国国家科学研究委员会早在1995年就指出，“研究者今天已经处在应用理论和计算来设计材料的初期阶段”。
经过数十年的发展，计算材料科学已涵盖了从微观的埃\textrm{(\AA)}到宏观的\textrm{(m)}量级的空间尺度和从飞秒\textrm{(fs)}到年\textrm{(Y)}的时间尺度，在高通量计算、数据挖掘及机器学习技术的加持下，形成了近年来推崇的``计算—数据—智能''驱动的新型研究范式。

%本文按照计算材料科学的发展历程，依次考察前计算机时代的理论奠基、快速发展阶段的方法勃兴、数据驱动的范式转型，以及跨尺度模拟中的核心软件与方法体系，最后对学科的机遇与挑战进行总结与展望。

%一、前计算机时代：理论基础与早期探索

%1.1 量子力学的奠基

计算材料科学的源头可上溯到上世纪初，随着量子力学的建立，\textrm{Bloch}将量子理论运用于周期体系，形成电子能带理论，\textrm{Hartree-Fock}应用平均场方法，近似求解分子的多电子体系波函数，开启了理论方法精确描述材料电子结构的时代。但是仅有理论不足以支持真实体系的有效描述。早在二十世纪二、三十年代交接之际，量子力学巨擘\textrm{P. A. M. Dirac}就敏锐地指出:``(描述)大部分物理和全部化学的基本定律的数学理论已然全部明了，唯一的困难在于这些方程是如此复杂，难以求解''\upcite{PRSLSA123-714_1929}。这句话点出了早期理论计算面临的最主要的困境——原理的自洽，在没有合适的工具时候，实际应用中依然一筹莫展。直到\textrm{1955}年， \textrm{C. W. Scherr}等用手摇计算机历时两年才完成\ch{N2}分子的\textrm{Hartree-Fock-Roothaan}从头计算，由此可见其艰难程度。%直到1960年，\textrm{L. Pauling}仍坦承：“自从薛定谔方程发现以来的30年中，我们看到，化学家感兴趣的物质性质只有很少几个作出了准确而又非经验性的量子力学计算”。这充分说明，尽管量子力学早已提供了描述材料行为的第一原理框架，但
在没有高性能计算机支撑的条件下，从头计算(第一原理)理论优势很难转化为实际的计算能力。

%的奠基工作。1913年，尼尔斯·玻尔建立了原子的量子模型，开启了人们对微观粒子运动规律的量化理解。20世纪20至30年代，量子力学实现了完整的理论建构，提出了一系列极具深远影响的基本原理。1928年，F. Bloch将量子理论运用于固体体系，奠定了固体电子能带理论的基调和后续研究的出发点。同年至1930年间，Hartree-Fock平均场方法的建立使得多电子体系的近似求解成为可能，为后续更具普遍性的密度泛函理论铺就了道路。1927年提出的Thomas-Fermi理论则开创了用电荷密度而非波函数来描述体系电子状态的思想先河。

%1.2 量子化学计算的艰苦起步

%然而，有了理论框架并不意味着实际计算就可以顺利进行。量子化学先驱P. A. M. Dirac在1930年便敏锐地指出：“解决全部化学的规律的数学方法已完全知道了，困难只是在于这些方程太复杂，无法求解”。这句话精确地概括了早期理论计算所面临的根本困境——原理上自洽，实践中却几乎寸步难行。


%1.3 固体物理与材料模拟的先驱性探索

%在计算能力极为有限的年代，固体物理学家和材料科学家仍然作出了一系列对后世影响深远的开创性探索。许多人开始在理论模型与近似假设之间寻找平衡，试图在不依赖大型计算设施的情况下触及材料微观世界的本质。例如，中国计算材料物理领域的奠基者王崇愚院士从含氧合金研究出发，提出“氧—层错复合体模型”等理论见解，这些早期工作在缺乏现代计算工具的条件下展现了深刻的理论洞察力。这些努力表明，即使在计算工具极为原始的环境中，材料科学家依然可以通过精巧的理论构建逼近复杂的材料行为。也可以说，正是在这个“有理论而无计算”的时期，人们最深刻地意识到：若要将理论付诸应用，必须等待电子计算机带来的范式革命。

%总而言之，前计算机时代奠定了计算材料学的两大支柱——量子力学理论体系与初步的实验性计算探索，但受限于计算工具的匮乏，这些理论成果多数停留在概念层级，未能充分转化为解决实际材料问题的系统工具。然而，正是这段漫长而艰难的积淀，为计算机时代到来的范式跃迁提供了不可或缺的理论起点。

%二、快速发展阶段：计算方法的大爆发

%2.1 计算机的兴起与摩尔定律的助推

二十世纪六十年代，电子计算机的诞生与飞速发展为计算材料科学赋予了全新的生命力。计算能力的跃迁带来了研究的实质性突破。正如现代计算材料学导论教材所示，从1946年至1971年，计算机经历了真空管、晶体管、集成电路直到大规模集成电路的数代更迭，%中央处理器\textrm{(CPU)}速度大致遵循摩尔定律所描述的每18个月翻一番的轨迹快速攀升。
这种指数级增长的算力资源使得过去需要手摇计算机耗时数年的体系，如今可以在可接受的时间尺度内完成解算，极大地释放了理论模型的潜能。

%2.2 密度泛函理论的创立与成熟

密度泛函理论（DFT）的提出无疑是计算材料学历史上最为关键的里程碑事件之一。尽管将体系基态能量表达为电子密度泛函的思想可以追溯到早期的Thomas-Fermi模型，但真正意义上的理论基石由Hohenberg和Kohn于1964年奠定。Hohenberg-Kohn第一定理指出，体系基态能量仅仅是电子密度的泛函，这一论断重新将复杂多电子问题的求解路径从高维波函数转向三维电子密度，为简化计算打开了全新通道。次年，Kohn与Sham进一步发表了KS方程，将实际电子体系映射为无相互作用的参考体系，从而使得DFT方法具备了可实际计算的形式。Jones通过综述研究指出，DFT计算的起源通常追溯到Hohenberg、Kohn和Sham在1964—1965年间发表的系列论文，而这些方法在化学和材料科学领域的广泛应用是在1990年之后才开始的。Kohn和Pople因该领域的卓越贡献在1998年获得诺贝尔化学奖，这一荣誉有力地确认了DFT在量子化学与材料计算领域的核心地位。

DFT的成功在于它巧妙地平衡了计算精度与计算成本。它不需要经验参数，仅以原子序数和基本物理常数为输入，就可以预测材料的电子结构、能带、磁性和弹性等诸多本征性质，为第一性原理材料模拟建立了坚实基础。

2.3 分子动力学的演进与Car-Parrinello方法

如果说DFT解决了“电子结构怎么算”的问题，那么分子动力学方法的发展则回答了“原子如何运动”的问题。1957—1959年，Alder和Wainwright首次提出并应用于理想“硬球”液体模型，开创了分子动力学模拟的先河。1963年，Rahman采用连续势模型对液态氩进行了分子动力学模拟，将模拟体系从理想化模型推进至实际材料体系。1967年，Verlet给出了著名的Verlet算法，这一简洁而高效的数值积分格式对分子动力学的长期稳定模拟至关重要。

1980年代初期，模拟条件控制技术取得重大进展：Anderson于1980年提出恒压分子动力学方法，同年Hoover对非平衡态分子动力学展开了系统研究，Parrinello和Rahman于1981年给出了恒定压强的分子动力学模型，Nose则在1984年提出了恒温分子动力学方法。这些进展使得分子动力学模拟日益接近实际的实验条件环境，其预测结果与实验数据的可比性由此大为提升。

然而，上述经典分子动力学方法的根本局限在于，原子间相互作用必须预先以经验力场或势函数的形式给出，而这对于含有化学键断裂与形成的复杂体系显然是不适用的。1985年，Car和Parrinello提出了革命性的第一性原理分子动力学方法，将电子结构的DFT计算与经典分子动力学框架有机结合，实现对电子和原子核运动的同时模拟。这一突破意味着，原子所受的力不再依赖预先给定的经验势函数，而是通过“实时”的量子力学计算获得，从而使得化学反应的动态过程研究首次进入原子模拟的射程之内。

2.4 蒙特卡罗方法与介观模拟方法的兴起

在多粒子体系的统计力学模拟领域，蒙特卡罗方法与分子动力学方法相辅相成。蒙特卡罗方法通过随机采样和马尔可夫过程来模拟系统的平衡态构型演化，在相图计算、自旋系统和材料生长动力学等方面具有独特优势，尤其擅长处理那些分子动力学难以直接涉及的状态空间高效采样问题。

与此同时，介观尺度模拟方法也在快速发展。介观尺度是连接原子/分子微观结构（纳米尺度）与宏观连续体（毫米/米尺度）的关键过渡尺度，其研究方法包括元胞自动机、相场动力学方法、位错动力学方法与多晶塑性有限元模型等。其中金兹堡—朗道型相场动力学方法通过引入连续相场变量描述界面演化，能够在介观尺度上模拟枝晶生长、晶粒粗化和畴壁运动等微结构形成过程，突破了微观模拟方法对空间和时间尺度的固有制约。

2.5 中国计算材料物理的奠基工作

在国际学术界快速发展的同时，中国的计算材料科学也在稳步推进。王崇愚院士是我国计算材料物理领域的奠基者与开拓者之一。他从含氧合金研究切入，提出了“氧—层错复合体模型”，并首创“多尺度物理参量解析传递算法”等核心算法，为多尺度模拟奠定了理论、方法与组织的多重基础。王崇愚还引入了离散变分方法，推动成立了中国金属学会离散变分分会、中国材料研究学会计算材料学分会，使我国的计算材料物理在大方向上有了组织化的平台与专门的学术交流渠道。他的成果曾应用于“两弹一星”等国家重点工程，在基础理论研究与重大工程需求之间架设了坚实的桥梁。

三、数据驱动的材料计算科学：从高通量计算到人工智能

3.1 材料基因组计划（MGI）与高通量计算的兴起

进入21世纪第二个十年，计算材料科学迎来了一次范式级的深刻变革。2011年，美国正式实施材料基因组计划，其核心理念是通过计算、数据和实验“三位一体”的方式，彻底变革传统以经验和实验为主导的“试错法”材料研发模式，将新材料从发现到开发、生产和应用的速度提升至传统模式的两倍。MGI强调建立“计算模拟和理论预测优先、实验验证在后”的新材料研发文化，这意味着材料发现的主导权正在从实验室逐步向计算平台迁移。

高通量计算是MGI的技术核心之一。“十三五”时期，我国自主开发了多种材料高通量、并发式、自动流程计算方法和软件，实现了万量级（10⁴级）高通量并发式计算，为高效、快速地筛选候选新材料和预测新物理效应奠定了坚实的基础，并实现了从建模、计算到数据分析的全程自动化、无人工干预的可视化操作。平台方面，国家超级计算天津中心依托我国自主研发的千万亿次（P级）和百亿亿次（E级）超级计算平台，构建了能够实现高通量材料计算和材料大数据应用管理的“中国材料高通量计算平台（CNMGE）”。

3.2 材料数据库的建设与发展

数据是数据驱动研究范式中最为核心的资源。在无机材料领域，先驱性工作如Materials Project的建立推动了覆盖整个元素周期表的计算材料数据库发展，其中包括AFLOW、OQMD、GNoME和OMat24数据集等。这些数据库集成了大规模第一性原理计算的结果，为全球材料研究人员提供了统一、开放和可复用的数据基础设施。

中国在这一领域同样取得了令人瞩目的进展。中国科学院物理研究所的Atomly材料数据库已将无机晶体材料扩充至30余万个，数据规模和质量已达到世界顶级水准，基于这一海量数据和高通量计算手段，科研团队在新型材料体系的快速搜索和预测方面取得了显著的成果。北京怀柔材料基因组研究平台还自研了具有自主知识产权的高通量作业流程序包，能够自动化调度高性能计算资源、分析计算结果并不断积累科学数据，构建了材料基因的核-心数据资源库。OPTIMADE API的提出进一步实现了AFLOW、Materials Cloud、Materials Project、NOMAD、OQMD等主要材料数据库的统一查询接口，使跨库数据检索与整合变得便捷易行。

3.3 机器学习与人工智能驱动的材料智能设计

近年来，机器学习和人工智能技术在材料科学领域的突破性应用，将数据驱动范式推向了一个崭新的高度。在材料科学领域，机器学习主要借助第一性原理数据支持新型化合物的大规模发现，并通过数字孪生等技术优化材料的制造工艺。监督式机器学习方法已被广泛应用于全球材料研发领域，其基本原理是通过构建实验数据数字化特征与目标性能在高维空间中的映射关系，实现对候选材料的快速评估和预测。

值得关注的是，数据驱动的材料研究已开始从单纯的性能预测向“材料创造”方向演进。用于材料发现的深度生成模型不再仅仅依赖已知数据中的显式关联，而是通过深度学习过程学习原子结构的隐式知识，并以这些知识为驱动实现新型材料结构的高效生成与创新设计。拓扑数据分析方法亦被引入材料构效关系研究，通过将材料微观结构映射为数学拓扑模型来实现科学、可解释的特征抽取。根知识与小数据协同的材料智能计算则提出了将系统论、动力学方程与化学还原物理涌现等深层次结构性知识融入模型训练的全新思路。中国自主开发的ALKEMIE平台整合了材料高通量自动计算模拟、材料数据库系统以及基于人工智能的数据挖掘技术，为数据驱动的材料研发新模式提供了面向实际需求的集成化工具支撑。

数据驱动范式的崛起并不意味着传统计算方法的退场。事实上，大规模材料数据库的价值恰恰在于它们内部的海量计算结果恰恰是以DFT、分子动力学等经典方法生成并积累的。机器学习模型从这些计算数据中学习到的规律，通过迁移学习等方式指导真实实验数据的预测，形成“计算预测—实验验证—数据反馈”的闭环优化流程，这正是下一代材料科学研究的核心图景所在。

四、材料模拟的软件与方法体系

计算材料科学领域最为鲜明的技术特征之一，是其“多尺度”属性。从电子波函数的量子尺度，到原子振动的微观尺度，再到晶粒和界面的介观尺度，直至工程构件的宏观尺度——每一个尺度对应的物理规律、模拟方法和计算工具都有显著差异。本节按微观尺度、介观尺度和宏观尺度三个层级的次序，对支撑这些尺度的核心方法及代表性软件进行系统介绍。

4.1 微观尺度

微观尺度通常指原子和分子层面（约0.1—10纳米），主要关注电子结构和原子排列对材料性质的决定性影响。这一尺度的核心计算方法是第一性原理密度泛函理论和量子化学方法。

密度泛函理论（DFT） 的计算流程通常包括：基于原子类型和初猜晶格常数构建初始结构模型；选择适当的交换关联泛函（如LDA、GGA等）；设定k点采样和截断能以控制计算精度；执行自洽场迭代求解Kohn-Sham方程，获得基态电子密度和总能量；进而计算能带结构、态密度、弹性常数、声子谱等物理属性。这一框架的不解自明之处在于：所有输出性质均从几个基本物理常数出发计算得出，无需任何经验输入，因此被称为“第一性原理”计算或“从头算（ab initio）”计算。

DFT方法的广泛适用性得益于大量高质量开源和商业软件的实现。VASP（Vienna Ab initio Simulation Package）是原子尺度材料建模领域的商业软件，使用投影缀加波方法和平面波基组，在固体物理和材料科学中近乎成为标准工具。Quantum ESPRESSO是集成的开源电子结构计算软件套件，支持平面波与赝势方法，涵盖基态计算、线性响应计算（用于声子色散曲线计算、介电常数和Born有效电荷）、电子−声子相互作用等广泛功能，遵循GNU公共许可证，可在高性能并行计算机到个人工作站上运行。此外，WIEN2k基于全势线性缀加平面波（FLAPW）方法，擅长处理强关联体系和磁性质。CASTEP也广泛应用于固体材料的第一性原理模拟。将每种软件的适用边界和特点联系起来看，VASP偏向更高精度和更广泛的元素覆盖，Quantum ESPRESSO以开源生态和功能全面见长，WIEN2k则在处理含d/f电子的过渡金属和稀土材料方面具有独特优势——研究团队可根据具体问题选择合适的工具。

分子动力学模拟——尽管在介观尺度层面同样重要——在微观尺度上同样至关重要，尤其是当模拟体系涉及数百至数十万原子、时间尺度达到皮秒至微秒量级时。经典的分子动力学通过积分Newton运动方程追踪原子轨迹，并从原子坐标和速度的统计分布中提取热力学量（温度、压力、扩散系数等）。其主要软件包括LAMMPS（大规模原子分子并行模拟器）、GROMACS（生物分子模拟主流工具）和NAMD等。

蒙特卡罗方法则从概率角度出发，通过重要性采样和Metropolis判据在构型空间中演化系统状态，尤其擅长平衡态统计力学计算和热力学积分，在合金短程有序、表面吸附构型搜索等场景中应用广泛。

4.2 介观尺度

介观尺度处于微观原子表象与宏观连续描述之间的过渡带（约10纳米至100微米），涉及晶粒、位错、第二相粒子、畴结构等微结构单元的集体行为。这一尺度的模拟方法在物理建模上比微观方法更粗粒化，但仍然保留了关键的结构信息。

相场方法是该尺度的代表性方法之一。相场方法引入连续相场变量来描述晶粒取向、成分或有序度，通过求解金兹堡—朗道型微分方程来模拟微结构随时间的演化，无需显式追踪复杂界面的几何形状。这一方法已被广泛用于金属凝固过程中的枝晶生长、固相相变中的组织粗化、铁电压电材料的畴转变以及锂离子插入时的相分离过程。

位错动力学直接从晶体位错的离散表现进行建模和演化模拟，关注位错线形核、滑移、交割及增殖过程，对理解金属塑性变形、加工硬化和疲劳行为具有重要价值。元胞自动机和多态波茨模型则为再结晶、晶粒生长等涉及拓扑重排的现象提供了易于并行实现且计算效率较高的模拟路径。

在软件工具层面，LAMMPS不仅支持经典分子动力学，还通过多种力场和改进的蒙特卡罗方法，被用于构建多元相图等复杂任务中的并行采样子模块。一些综合性计算平台将介观模拟与微观计算进行集成，通过网格接口调度异构计算资源，将从分子动力学中获取的参数传递给相场模拟使用，从而实现跨尺度的数据链贯通。

将微观尺度和介观尺度的关联加以审视便不难理解：“微观输入，介观响应”正是多尺度模拟的核心逻辑。DFT输出弹性常数和层错能，分子动力学输出扩散系数和界面能——这些经参数化处理后成为介观模拟的“数据底层”，使得介观尺度的预测建立在原子尺度物理规律的基础之上，而非纯粹的唯象模型。

4.3 宏观尺度

宏观尺度包括米量级工程构件乃至更大尺度的结构分析，其核心方法是有限元方法（FEM） 。有限元方法将连续体离散化为有限数量的单元，在每个单元上求解弹塑性力学、传热学、流体力学等偏微分方程，最终通过整体矩阵组装获得构型全场的力学响应。

在工程应用层面，宏观尺度模拟几乎无法离开商业或通用有限元软件。ABAQUS被广泛用于非线性结构力学、热分析和复合材料建模等领域，支持晶体塑性有限元方法的用户材料子程序开发和内外多尺度耦合计算。COMSOL Multiphysics擅长多物理场耦合问题（电磁—热—力—流等），其用户友好的图形化界面和内置物理场接口使其在多尺度建模中广受欢迎。ANSYS和Fluent则是航空航天与能源工程领域的通用仿真利器。

在计算方法论层面，宏观尺度模拟面临的根本问题是本构关系的确定。工程材料在宏观尺度下的行为本质上是其微观结构统计平均的结果。以纤维增强复合材料为例，其宏观弹性模量和破坏强度依赖于纤维取向分布、界面结合强度以及基体中的孔隙率等多重微结构因素。如何在宏观有限元模型中有效反映这些微观特征——正是多尺度方法的核心课题。现代解决思路包括：FE²方法通过在宏观积分点上嵌入微观有限元模型来实时传递本构信息、均匀化方法通过统计平均将微结构响应上升为宏观本构参数、以及准连续介质方法在关键局部区域保留原子尺度细节并采取粗粒化近似加以衔接。国家相关部门开发的高通量材料集成计算平台支持与高性能计算集群的动态绑定，实现了从微观建模到计算再到数据分析的全程自动化与可视化操作，为宏观多尺度耦合提供了坚实的平台支撑。

4.4 跨尺度协同：从贯通到融合

计算材料科学在尺度上的完整性并非仅仅依赖于各尺度独立模拟工具的加和，而是要建立多种模拟方法之间的数据流与参数传递通道。

在层级耦合方面，微观尺度的DFT计算导出特定材料的弹性常数、扩散势垒和表面能等基础参数，这些参数被力场参数化后用于分子动力学模拟；分子动力学输出的中间尺度性质进一步被粗粒化，输入相场模型以模拟微米至毫米级的微结构演化；最终，相场模拟给出的统计性信息被嵌入宏观有限元计算的本构关系中。这一链条在空间尺度上跨越了约10个数量级，在时间尺度上跨越了约15个数量级，其内在挑战不言自明。

在并发耦合方面，一些前沿工作已尝试在单个模拟中将不同尺度的求解耦合在同一时间步中进行。例如，分子动力学和有限元耦合可以在裂尖附近原子区采用分子动力学、在其外围连续区采用有限元方法，从而在不丢失关键区域物理细节的前提下大幅扩展整体模拟规模。准连续介质方法正是这一思路的一种系统化实现，它通过追踪代表性原子并对其余原子的运动进行插值近似，实现了模拟位错运动、裂纹扩展等复杂现象时计算效率与精度的合理平衡。

从更长远的发展趋势看，近年来机器学习和材料大数据方法的加入了为多尺度耦合提供了新的解决方案。例如，从高通量DFT计算数据中训练原子间机器学习势函数（如等变图神经网络势），可以使分子动力学模拟在保持量子力学级精度的同时，将可模拟的原子规模从几千提升至百万甚至十亿量级，从而消弭微观模拟与原子统计之间的尺度鸿沟。这一趋势预示着“数据驱动”与“多尺度建模”两条路径正不断交汇，共同推动计算材料科学从“经验性校正”迈向“预测性设计”。

五、机遇与挑战

计算材料学已从理论物理的一个分支成长为驱动新材料研发的核心技术引擎，但它仍面临一系列亟待深入的关键问题。

第一是多尺度物理信息的贯通问题。DFT面对包含数千以上原子的体系时计算成本过高，而机器学习势函数虽能显著加速模拟，但其训练数据覆盖范围有限，对体系化学组成和结构相空间的泛化能力仍存在不确定性，如何在不同精度与不同效率的尺度模型之间建立无缝的贯通机制是当前研究的核心瓶颈之一。

第二是高通量计算精度与数据质量的可信度问题。不同计算代码、不同交换关联泛函、不同赝势和收敛参数设置均会给DFT计算结果带来不可忽视的差异。对AFLOW、Materials Project和OQMD三大数据库进行对比的研究已经揭示了材料性质预测在不同数据库之间的一致性问题。如何统一计算标准、量化计算不确定性并与实验数据同化和适配，是材料基因组计划走向成熟应用的临界条件。

第三是数据驱动与理论驱动之间的有机融合难题。单纯依赖机器学习从大数据中挖掘规律，固然可在某些体系上取得令人满意的预测效果，但如果不注入物理约束，模型可能产生非物理输出，在训练数据分布之外区域会表现出极大的预测不确定性。物理学神经网络将守恒律、对称性等先验知识嵌入了模型架构，以根知识与小数据为核心的协同计算架构也正在探索如何将系统论、动力学方程与非线性数据结构进行整合。

第四是高通量材料计算平台的国产化与自主化迫切性。尽管我国在CNMGE平台和ALKEMIE融合平台方面取得了进展，但从核心计算软件到数据库体系都仍面临与国外先进体系有效衔接和竞争的现实压力。我国自主开发的Hylanemos等第一性原理软件在计算性能与精度方面已达到国际先进水平，但在生态完整度、用户人群规模和持续维护机制上仍有发展空间。

本文参考了王崇愚院士的学术工作与相关领域的前沿发展。2026年3月王崇愚院士在北京逝世，为中国计算材料物理界带来了巨大的损失。他为中国计算材料物理学科的奠基、开拓和多尺度方法的创新所做的不朽贡献，以及他对“材料基因组”计划在国内实施的大力推动，将持续深远地影响整个学科的发展方向。谨以此表达深切缅怀。

参考文献

[1] 计算材料学. 百度百科.（引用于学科定义及核心方法概述）

[2] 杨小渝, 任杰, 王娟, 等. 基于材料基因组计划的计算和数据方法[J]. 科技导报, 2016,34(24):62-67.（引用于材料基因组计划的核心理念与高通量计算平台）

[3] Jones R O. Density functional theory: Its origins, rise to prominence, and future[J]. Reviews of Modern Physics, 2015,87:897.（引用于DFT起源与密度泛函方法发展）

[4] R. O. Jones. Density functional theory: Its origins, rise to prominence, and future[J]. Rev. Mod. Phys., 2015,87:897-923.（引用于DFT理论基础及发展）

[5] Alder B J, Wainwright T E. Phase transition in elastic disks[J]. Phys. Rev., 1957.（引用于分子动力学起源）

[6] Car R, Parrinello M. Unified approach for molecular dynamics and density-functional theory[J]. Phys. Rev. Lett., 1985,55:2471-2474.（引用于第一性原理分子动力学方法）

[7] Parrinello M, Rahman A. Polymorphic transitions in single crystals: A new molecular dynamics method[J]. J. Appl. Phys., 1981,52:7182-7190.（引用于恒压分子动力学方法）

[8] Hohenberg P, Kohn W. Inhomogeneous electron gas[J]. Phys. Rev., 1964,136:B864-B871.（引用于密度泛函理论的Hohenberg-Kohn定理）

[9] Kohn W, Sham L J. Self-consistent equations including exchange and correlation effects[J]. Phys. Rev., 1965,140:A1133-A1138.（引用于Kohn-Sham方程）

[10] 王崇愚院士生平与贡献. 光明日报, 2026年3月.（引用于王崇愚院士的奠基性工作与学术贡献）

[11] Raabe D. 计算材料学[M]. 北京:化学工业出版社, 2003.（引用于介观尺度模拟方法与尺度分级论述）

[12] 中国材料高通量计算平台建设进展. 国家超级计算天津中心. 2020-2025.（引用于高通量计算平台与ALKEMIE集成）

[13] A new horizon in data-driven materials research. EurekAlert, 2025.（引用于材料数据库与数据驱动材料研究趋势）

[14] Atomly材料数据库. 中国科学院物理研究所.（引用于Atomly数据库规模与自主研发成果）

[15] OPTIMADE API for Materials Databases. 中科院物理所报告, 2023.（引用于主要材料数据库统一搜索查询机制）
